\section*{ЗАКЛЮЧЕНИЕ}
В данной работе было проведено численное исследование модели $G(n,p)$. Были экспериментально подтверждены известные теоретические факты 
о размерах компонент связности, о связности графа, о длине пути.

Теорема о размере компоненты связности при $p = \frac{1+\epsilon}{n}$ выполнялась при $n=10000$, $\epsilon=0.1$ в 
94\% случаях, при $\epsilon = 0.5$ и $\epsilon=0.01$  
выполнялась в 100 \% случаях.

Оба фазовых перехода модели $G(n,p)$ наблюдались в численных экспериментах. После достижения  $p = \frac{1}{n}$ 
в графе появлялась гигантская компонента связности, размер
которой рос достаточно быстро. После достижения 
$p = \frac{\ln{n}}{n}$ доля связных графов среди всех экспериментов достаточно быстро росла.

Тест на симметричность распределения размера наибольших
компонент связности показал отсутствие симметричности 
почти при всех значениях $p$. Такой же результат дал
тест на симметричность распределения длин наибольших найденных путей.

Фазовые переходы также заметны при исследовании структуры графов. 
До первого фазового перехода $p=\frac{1}{n}$ 
графы представляют собой множество деревьев, после 
его достижения начинают появляться подграфы с циклами, 
после второго фазового перехода $p = \frac{\ln{n}}{n}$
из деревьев остаются немногочисленные изолированные вершины.

Таким образом, численное исследование подтвердило теоретические 
свойства модели $G(n,p)$ и продемонстрировало 
существование фазовых переходов. Перспективой дальнейшего исследования является изучение
поведения других моделей случайных графов, таких 
как модель Барабаши-Альберт и модель малого мира Уоттса-Строгаца.
\addcontentsline{toc}{section}{ЗАКЛЮЧЕНИЕ}
